\chapter{A Script-Unified Subword Vocabulary}
\label{ch:method}

Telugu, Kannada and Tamil are related languages written in different scripts. We transliterate
all three into a common representation before learning subwords, so that cognates share units,
and transliterate the output back afterwards.

\begin{figure}[htbp]
  \centering
  \begin{tikzpicture}[node distance=5mm and 7mm,
      box/.style={draw=accent, rounded corners=2pt, minimum height=9mm, align=center, font=\small}]
    \node[box] (s) {Source text\\(native script)};
    \node[box, right=of s] (t) {Common\\representation};
    \node[box, right=of t] (m) {Multilingual\\model};
    \node[box, right=of m] (o) {English\\output};
    \draw[-{Stealth}] (s) -- (t); \draw[-{Stealth}] (t) -- (m); \draw[-{Stealth}] (m) -- (o);
  \end{tikzpicture}
  \caption{Script unification before subword segmentation.}
  \label{fig:unify}
\end{figure}
